\section{Method}
% Describe the theoretical framework, experiment, data and evaluation.

\paragraph{Data and release.} We use the training split of the UCI Adult census dataset
% AUDIT: the historical cleaned-row count needs the original cache; the current cache has 32,561 complete rows.
  \citep{becker1996adult}. To keep the search for trees and the reconstruction tractable, all
  features are binary: age and hours worked per week are split at their medians, capital
  gain into zero or positive, marital status into four groups and education into five (one
  column per group), and sex is a single column. This gives 13 columns. The label $Y$ indicates
  whether a person earns more than 50K; the attacker tries to reconstruct sex. Keeping the same binary target across conditions
  makes the reconstruction errors directly comparable. For each seed
  we draw a private training sample of $n=100$ rows and a separate public sample of 2,000
  other rows. On the private sample, TreeFARMS \citep{xin2022treefarms} (regularization
  0.01) lists every tree whose regularized objective is within 10\% of the best one: the
  Rashomon set. The publisher releases $m$ of these trees together with the exact number of
  training rows of each class in every leaf, which we call the \emph{release}. Trees are
  added either in a random order (random selection) or in the order in which TreeFARMS lists
  them (sequential selection), so a larger release always contains the smaller ones. The attacker knows every row's income label.

\paragraph{Reconstruction attack (SQ1, SQ2).} We adapt the constraint programming model of
  DRAFT \citep{ferry2024draft} to trees trained on the whole sample, without bootstrap
  sampling. The unknown data form a table of 100 rows and 13 binary columns. Rows must follow paths consistent with their values in every tree, match all released
  leaf class counts, and have exactly one active column per one-hot group. The solver
  CP-SAT\footnote{OR-Tools 9.15, \url{https://developers.google.com/optimization}.} searches
  for any table that meets these constraints (180\,s and 4 workers per solve). Such a table is
  consistent with the release, but it does not have to be the original data.

  For SQ1 we compare four kinds of background knowledge: none, A, B and A+B, where A is
  marital status and B is education in one run and age in another. This compares binary and grouped categorical knowledge, alone and alongside
  marital status, to test whether the number of known attributes captures their benefit. Known attributes are fixed to their true values for every training row. Both runs use seeds 0 to 99 and 14 release sizes from 1 to 200, skipping sizes
  larger than the Rashomon set. For SQ2 we repeat A+B with sequential selection.

  We score anonymous rows by matching reconstructed and true rows within each income
  group, minimizing Hamming distance over non-sex values, then counting sex errors. Only the evaluator can do this pairing, because it uses the true data. As controls that use neither the trees nor A or
  B, we guess sex uniformly at random, or according to its frequency in the public sample.
% AUDIT: the majority baseline needs the original private-sample cache; it is not saved in the CSVs.

\paragraph{Record-level confidence (SQ3).} Confidence must use only attacker information to assess individual claims.
  This separate experiment uses the same preprocessing and constraints, no background
  knowledge, and full records (13 columns plus income). The attacker outputs 100 rows
  respecting released income totals.
  The $k$-th copy of record $r$ claims at least $k$ occurrences in the private sample.
  The evaluator checks counts without identifying people. Ideally, among claims assigned 80\%
  confidence, about 80\% would be correct, without fitting scores to private ground truth.

\paragraph{Prior-guided sampling.} The attacker draws $K=20$ reference tables from the
  public sample, sampling each row with the same income label. For each reference table,
  CP-SAT minimizes the number of changed cells subject to release consistency (30\,s and
  6 workers per solve); the solver may return a table without proving optimality. A claim's confidence
  is the share of returned tables containing at least $k$ copies of $r$. For each income
  label, the attacker selects the highest-confidence claims. Private data are used only
  for evaluation. Confidence measures agreement across the returned reconstructions, not certainty
  that a claim is correct.

\paragraph{Controls and variants.} As a control without trees, the attacker estimates the
  probability of at least $k$ copies of $r$ from the public sample alone, with a binomial
  model and additive smoothing of 0.1 (``public statistics only''). With only 200 public records, repeated sampling may overstate agreement.
  The Dirichlet variant first draws a distribution over the 640 possible records using
  public counts plus 0.1, then samples each reference table from it; its no-tree control
  uses a beta-binomial model. This accounts for uncertainty in public record frequencies,
  rather than treating the small public sample as the population. All confidence
  runs use random selection, no background knowledge, 20 seeds and $m\in\{10,25,50,100\}$.
  We report accuracy (the share of correct claims), overconfidence (mean confidence minus
  accuracy), AUROC, and how many claims reach a confidence of at least 0.9 and how many of
  those are correct.

\paragraph{Statistics.} Each seed is one paired observation, because all conditions in a
  comparison share its samples and Rashomon set. We use the two-sided Wilcoxon signed-rank
  test \citep{wilcoxon1945individual}. For SQ1, each seed's gain averages the error reduction against no background knowledge
  over shared release sizes; Table~\ref{tab:attack} reports the mean across seeds. These six tests
  (three conditions in two runs) are corrected with Holm's method \citep{holm1979multiple}.
  For SQ3 we compare accuracy across paired seeds, or the absolute gap between mean
  confidence and accuracy for each seed.
